% GENERATED FILE. Edit canonical lessons, then run npm run generate:books.
% canonical-source-hash: fnv1a64:38e579300c2b285f
% canonical-lessons: HI-C340-upanyas, HI-C340-nayak, HI-C340-patra, HI-C340-adhyay, HI-C340-anjam

\chapter{A Novel}
\label{ch:hi-a-novel}
\hlchaptermodality{340}

\begin{chapteropening}
\textbf{By the end of this chapter:} \emph{I can say what a novel is about and who is in it.}

Complete the last lesson of chapter 340: I can say what a novel is about and who is in it.
\end{chapteropening}

\section[\texorpdfstring{\hi{उपन्यास} (upanyās)}{उपन्यास (upanyās)}]{\hi{उपन्यास} (upanyās) — a novel}

\label{lesson:HI-C340-upanyas}



\begin{quote}
Before the new one: say the Hindi for trickery, then the Hindi for to lure.
\end{quote}

\subsection*{You'll want to know: \hi{उपन्यास}}
\textbf{\hi{उपन्यास}} — \emph{upanyās} — \textquotedblleft{}a novel\textquotedblright{}.

\begin{quote}
\hi{मीरा} \hi{ने} \hi{छुट्टी} \hi{में} \hi{एक} \hi{लंबा} \hi{उपन्यास} \hi{पढ़ा।}
\end{quote}

\begin{quote}
\emph{Mīrā ne chuṭṭī meṁ ek lambā upanyās paṛhā.}
\end{quote}

\begin{quote}
Mira read a long novel during the holidays.
\end{quote}

Masculine.

\begin{cousinweb}
\textbf{\hi{उपन्यास}} (\emph{upanyās}) is Sanskrit \emph{upanyāsa}, \textquotedblleft{}a placing near; a statement\textquotedblright{} (Monier-Williams). Modern Hindi uses it for a novel; the dictionary gives only the older senses.
\end{cousinweb}

\begin{grammarlens}[title={What you've built}]
One more word for telling what happened.
\end{grammarlens}

\subsection*{Your turn}
\begin{itemize}
\raggedright
  \item \emph{Say it:} \emph{upanyās}
  \item \emph{Say it:} \emph{upanyās}, once more
  \item \emph{Say it:} \emph{Mīrā ne chuṭṭī meṁ ek lambā upanyās paṛhā.} — Mira read a long novel during the holidays.
  \item \emph{Recall:} say the Hindi for trickery, then the Hindi for to lure, then say \emph{upanyās} again
\end{itemize}

\begin{tcolorbox}[breakable,colback=teal!4,colframe=teal!35!black,title={Before you move on}]
What does \hi{उपन्यास} mean? (\textquotedblleft{}A novel\textquotedblright{}.) Say it once more.
\end{tcolorbox}
\section[\texorpdfstring{\hi{नायक} (nāyak)}{नायक (nāyak)}]{\hi{नायक} (nāyak) — a hero (of a story); a leader}

\label{lesson:HI-C340-nayak}



\begin{quote}
Before the new one: say the Hindi for to lure, then the Hindi for a novel.
\end{quote}

\subsection*{You'll want to know: \hi{नायक}}
\textbf{\hi{नायक}} — \emph{nāyak} — \textquotedblleft{}a hero (of a story); a leader\textquotedblright{}.

\begin{quote}
\hi{उपन्यास} \hi{का} \hi{नायक} \hi{एक} \hi{ग़रीब} \hi{किसान} \hi{है।}
\end{quote}

\begin{quote}
\emph{upanyās kā nāyak ek ġarīb kisān hai.}
\end{quote}

\begin{quote}
The hero of the novel is a poor farmer.
\end{quote}

Masculine.

\begin{cousinweb}
\textbf{\hi{नायक}} (\emph{nāyak}) is Sanskrit \emph{nāyaka}, \textquotedblleft{}a leader\textquotedblright{}, and already in Sanskrit drama the hero (Monier-Williams).
\end{cousinweb}

\begin{grammarlens}[title={What you've built}]
One more word for telling what happened.
\end{grammarlens}

\subsection*{Your turn}
\begin{itemize}
\raggedright
  \item \emph{Say it:} \emph{nāyak}
  \item \emph{Say it:} \emph{nāyak}, once more
  \item \emph{Say it:} \emph{upanyās kā nāyak ek ġarīb kisān hai.} — The hero of the novel is a poor farmer.
  \item \emph{Recall:} say the Hindi for to lure, then the Hindi for a novel, then say \emph{nāyak} again
\end{itemize}

\begin{tcolorbox}[breakable,colback=teal!4,colframe=teal!35!black,title={Before you move on}]
What does \hi{नायक} mean? (\textquotedblleft{}A hero (of a story); a leader\textquotedblright{}.) Say it once more.
\end{tcolorbox}
\section[\texorpdfstring{\hi{पात्र} (pātra)}{पात्र (pātra)}]{\hi{पात्र} (pātra) — a character (in a story)}

\label{lesson:HI-C340-patra}



\begin{quote}
Before the new one: say the Hindi for a novel, then the Hindi for a hero (of a story).
\end{quote}

\subsection*{You'll want to know: \hi{पात्र}}
\textbf{\hi{पात्र}} — \emph{pātra} — \textquotedblleft{}a character (in a story)\textquotedblright{}.

\begin{quote}
\hi{इस} \hi{कहानी} \hi{में} \hi{सिर्फ़} \hi{तीन} \hi{पात्र} \hi{हैं।}
\end{quote}

\begin{quote}
\emph{is kahānī meṁ sirf tīn pātra haiṁ.}
\end{quote}

\begin{quote}
There are only three characters in this story.
\end{quote}

Masculine.

\begin{cousinweb}
\textbf{\hi{पात्र}} (\emph{pātra}) is Sanskrit \emph{pātra} (Monier-Williams), first a vessel and then a worthy person and an actor in a play; in Hindi it is any character in a story.
\end{cousinweb}

\begin{grammarlens}[title={What you've built}]
One more word for telling what happened.
\end{grammarlens}

\subsection*{Your turn}
\begin{itemize}
\raggedright
  \item \emph{Say it:} \emph{pātra}
  \item \emph{Say it:} \emph{pātra}, once more
  \item \emph{Say it:} \emph{is kahānī meṁ sirf tīn pātra haiṁ.} — There are only three characters in this story.
  \item \emph{Recall:} say the Hindi for a novel, then the Hindi for a hero (of a story), then say \emph{pātra} again
\end{itemize}

\begin{tcolorbox}[breakable,colback=teal!4,colframe=teal!35!black,title={Before you move on}]
What does \hi{पात्र} mean? (\textquotedblleft{}A character (in a story)\textquotedblright{}.) Say it once more.
\end{tcolorbox}
\section[\texorpdfstring{\hi{अध्याय} (adhyāy)}{अध्याय (adhyāy)}]{\hi{अध्याय} (adhyāy) — a chapter of a book}

\label{lesson:HI-C340-adhyay}



\begin{quote}
Before the new one: say the Hindi for a hero (of a story), then the Hindi for a character (in a story).
\end{quote}

\subsection*{You'll want to know: \hi{अध्याय}}
\textbf{\hi{अध्याय}} — \emph{adhyāy} — \textquotedblleft{}a chapter of a book\textquotedblright{}.

\begin{quote}
\hi{पहले} \hi{अध्याय} \hi{में} \hi{नायक} \hi{गाँव} \hi{छोड़ता} \hi{है।}
\end{quote}

\begin{quote}
\emph{pahle adhyāy meṁ nāyak gā̃v choṛtā hai.}
\end{quote}

\begin{quote}
In the first chapter the hero leaves his village.
\end{quote}

Masculine.

\begin{cousinweb}
\textbf{\hi{अध्याय}} (\emph{adhyāy}) is Sanskrit \emph{adhyāya}, \textquotedblleft{}a reading, a lesson\textquotedblright{}, and so a chapter (Monier-Williams).
\end{cousinweb}

\begin{grammarlens}[title={What you've built}]
One more word for telling what happened.
\end{grammarlens}

\subsection*{Your turn}
\begin{itemize}
\raggedright
  \item \emph{Say it:} \emph{adhyāy}
  \item \emph{Say it:} \emph{adhyāy}, once more
  \item \emph{Say it:} \emph{pahle adhyāy meṁ nāyak gā̃v choṛtā hai.} — In the first chapter the hero leaves his village.
  \item \emph{Recall:} say the Hindi for a hero (of a story), then the Hindi for a character (in a story), then say \emph{adhyāy} again
\end{itemize}

\begin{tcolorbox}[breakable,colback=teal!4,colframe=teal!35!black,title={Before you move on}]
What does \hi{अध्याय} mean? (\textquotedblleft{}A chapter of a book\textquotedblright{}.) Say it once more.
\end{tcolorbox}
\section[\texorpdfstring{\hi{अंजाम} (añjām)}{अंजाम (añjām)}]{\hi{अंजाम} (añjām) — an outcome; the end}

\label{lesson:HI-C340-anjam}



\begin{quote}
Before the new one: say the Hindi for a character (in a story), then the Hindi for a chapter of a book.
\end{quote}

\subsection*{You'll want to know: \hi{अंजाम}}
\textbf{\hi{अंजाम}} — \emph{añjām} — \textquotedblleft{}an outcome; the end\textquotedblright{}.

\begin{quote}
\hi{उपन्यास} \hi{का} \hi{अंजाम} \hi{मुझे} \hi{पसंद} \hi{नहीं} \hi{आया।}
\end{quote}

\begin{quote}
\emph{upanyās kā añjām mujhe pasand nahīṃ āyā.}
\end{quote}

\begin{quote}
I didn't like the way the novel ended.
\end{quote}

Masculine.

\begin{cousinweb}
\textbf{\hi{अंजाम}} (\emph{añjām}) is usually traced to Persian \emph{anjām}, \textquotedblleft{}an end, a completion\textquotedblright{}. It is how things turn out: \textbf{\hi{इसका} \hi{अंजाम} \hi{बुरा} \hi{होगा}}, this will end badly.
\end{cousinweb}

\begin{grammarlens}[title={What you've built}]
One more word for telling what happened.
\end{grammarlens}

\subsection*{Your turn}
\begin{itemize}
\raggedright
  \item \emph{Say it:} \emph{añjām}
  \item \emph{Say it:} \emph{añjām}, once more
  \item \emph{Say it:} \emph{upanyās kā añjām mujhe pasand nahīṃ āyā.} — I didn't like the way the novel ended.
  \item \emph{Recall:} say the Hindi for a character (in a story), then the Hindi for a chapter of a book, then say \emph{añjām} again
\end{itemize}

\begin{tcolorbox}[breakable,colback=teal!4,colframe=teal!35!black,title={Before you move on}]
What does \hi{अंजाम} mean? (\textquotedblleft{}An outcome; the end\textquotedblright{}.) Say it once more.
\end{tcolorbox}
